\documentclass[11pt]{article}
\usepackage[margin=1in]{geometry}
\usepackage{xcolor}
\usepackage{tcolorbox}
\usepackage{hyperref}
\usepackage{enumitem}
\setlist{noitemsep,topsep=0pt}
\usepackage{booktabs}
\usepackage{array}
\usepackage{amsmath}
\usepackage{graphicx}

% Define OSU orange and AI blue colors
\definecolor{osaorange}{RGB}{220,115,38}
\definecolor{aiblue}{RGB}{30,144,255}

% Configure hyperref colors
\hypersetup{
    colorlinks=true,
    linkcolor=blue,
    urlcolor=blue,
    citecolor=blue
}

\title{}
\author{}
\date{}

\begin{document}

% Orange header box with black outline
\begin{tcolorbox}[colback=osaorange,colframe=black,boxrule=2pt,arc=0pt,left=6pt,right=6pt,top=10pt,bottom=10pt]
\centering
\textcolor{white}{\textbf{\Large Week 9 In-Class Worksheet}}\\
\textcolor{white}{\textbf{\large Correlation and Simple Linear Regression}}
\end{tcolorbox}

\vspace{8pt}

\noindent\textbf{Name:} \rule{8cm}{0.4pt} \hfill \textbf{Date:} \rule{4cm}{0.4pt}\\[0.5ex]
\textbf{Course:} H524 Introduction to Biostatistics\\
\textbf{Topic:} Correlation, Regression, and Prediction

\vspace{8pt}

\section*{Instructions}
Work with a partner to solve the following problems. Show all your work and calculations. You may use R to verify answers, but demonstrate your understanding by showing the steps.

\vspace{12pt}

\section{Problem 1: Correlation Analysis}

A researcher studies the relationship between daily calcium intake (mg) and bone density (g/cm$^2$) in 8 postmenopausal women:

\vspace{0.3cm}

\begin{center}
\begin{tabular}{ccc}
\toprule
Subject & Calcium (mg) & Bone Density (g/cm$^2$) \\
\midrule
1 & 800 & 0.85 \\
2 & 1000 & 0.90 \\
3 & 1200 & 0.95 \\
4 & 900 & 0.88 \\
5 & 1100 & 0.92 \\
6 & 1300 & 0.98 \\
7 & 950 & 0.89 \\
8 & 1050 & 0.91 \\
\bottomrule
\end{tabular}
\end{center}

\vspace{0.5cm}

\textbf{Part A:} Create a scatter plot of bone density (Y) vs. calcium intake (X). Describe the relationship you observe.

\vspace{4cm}

\textbf{Part B:} Calculate summary statistics:

\vspace{0.3cm}

Mean calcium intake: $\bar{X} = $ \rule{3cm}{0.4pt}

\vspace{0.3cm}

Mean bone density: $\bar{Y} = $ \rule{3cm}{0.4pt}

\vspace{0.3cm}

SD of calcium: $s_X = $ \rule{3cm}{0.4pt}

\vspace{0.3cm}

SD of bone density: $s_Y = $ \rule{3cm}{0.4pt}

\vspace{1cm}

\textbf{Part C:} The Pearson correlation coefficient is $r = 0.94$. Interpret this value in context.

\vspace{3cm}

\newpage

\textbf{Part D:} Test whether the correlation is statistically significant at $\alpha = 0.05$.

\vspace{0.3cm}

Hypotheses:

$H_0$: \rule{5cm}{0.4pt}

$H_A$: \rule{5cm}{0.4pt}

\vspace{0.5cm}

Test statistic: $t = \frac{r\sqrt{n-2}}{\sqrt{1-r^2}}$

\vspace{3cm}

Degrees of freedom: $df = $ \rule{2cm}{0.4pt}

\vspace{0.5cm}

Critical value (from t-table): $t_{0.975, df} = $ \rule{2cm}{0.4pt}

\vspace{0.5cm}

Decision: \rule{8cm}{0.4pt}

\vspace{0.5cm}

Conclusion (in context):

\vspace{2cm}

\textbf{Part E:} Calculate the coefficient of determination, $r^2$. Interpret this value.

\vspace{0.3cm}

$r^2 = $ \rule{3cm}{0.4pt}

\vspace{0.3cm}

Interpretation:

\vspace{2cm}

\newpage

\section{Problem 2: Simple Linear Regression}

Using the same data from Problem 1, we want to predict bone density from calcium intake.

\vspace{0.3cm}

\textbf{Part A:} Calculate the regression slope using the formula:

$$b_1 = r \times \frac{s_Y}{s_X}$$

Given: $r = 0.94$, $s_Y = 0.042$ g/cm$^2$, $s_X = 167.7$ mg

\vspace{3cm}

\textbf{Part B:} Calculate the regression intercept using:

$$b_0 = \bar{Y} - b_1 \bar{X}$$

Given: $\bar{Y} = 0.910$ g/cm$^2$, $\bar{X} = 1037.5$ mg

\vspace{3cm}

\textbf{Part C:} Write the regression equation:

$$\widehat{\text{Bone Density}} = $$ \rule{8cm}{0.4pt}

\vspace{1cm}

\textbf{Part D:} Interpret the slope in context. What does it mean biologically?

\vspace{3cm}

\textbf{Part E:} Is the intercept meaningful in this context? Why or why not?

\vspace{2.5cm}

\newpage

\section{Problem 3: Predictions and Residuals}

Continue with the regression model from Problem 2.

\vspace{0.3cm}

\textbf{Part A:} Use your regression equation to predict bone density for a woman with calcium intake of 1100 mg/day.

\vspace{0.3cm}

Calculation:

\vspace{3cm}

Predicted bone density: \rule{4cm}{0.4pt}

\vspace{1cm}

\textbf{Part B:} Subject 5 from the original data has calcium intake 1100 mg and observed bone density 0.92 g/cm$^2$. Calculate the residual for this subject.

\vspace{0.3cm}

Residual = Observed $-$ Predicted = \rule{6cm}{0.4pt}

\vspace{3cm}

\textbf{Part C:} Is this residual positive or negative? What does that tell you about this subject?

\vspace{3cm}

\textbf{Part D:} Would you feel confident predicting bone density for someone with calcium intake of 2000 mg/day? Why or why not?

\vspace{3cm}

\newpage

\section{Problem 4: Blood Pressure and Age}

A study of 12 adults examines the relationship between age and systolic blood pressure. The regression output is:

\vspace{0.3cm}

\begin{center}
\begin{tabular}{lrrrr}
\toprule
& Estimate & Std. Error & t value & $p$-value \\
\midrule
Intercept & 80.5 & 8.2 & 9.82 & $<0.001$ \\
Age & 1.15 & 0.18 & 6.39 & $<0.001$ \\
\bottomrule
\end{tabular}
\end{center}

Residual standard error: 6.3 mmHg on 10 degrees of freedom\\
Multiple R-squared: 0.803, Adjusted R-squared: 0.783

\vspace{0.5cm}

\textbf{Part A:} Write the regression equation.

\vspace{2cm}

\textbf{Part B:} Interpret the slope coefficient in the context of this study.

\vspace{3cm}

\textbf{Part C:} Test whether there is a significant relationship between age and blood pressure at $\alpha = 0.05$.

\vspace{0.3cm}

Hypotheses:

$H_0$: \rule{5cm}{0.4pt}

$H_A$: \rule{5cm}{0.4pt}

\vspace{0.5cm}

Test statistic: $t = $ \rule{3cm}{0.4pt} (from table above)

\vspace{0.5cm}

$p$-value: \rule{3cm}{0.4pt}

\vspace{0.5cm}

Decision and conclusion:

\vspace{2.5cm}

\newpage

\textbf{Part D:} Construct a 95\% confidence interval for the slope $\beta_1$.

\vspace{0.3cm}

Formula: $b_1 \pm t_{0.975, 10} \times \text{SE}(b_1)$

\vspace{0.3cm}

Given: $t_{0.975, 10} = 2.228$, $\text{SE}(b_1) = 0.18$

\vspace{3cm}

95\% CI: \rule{6cm}{0.4pt}

\vspace{0.5cm}

Interpretation (in context):

\vspace{2.5cm}

\textbf{Part E:} What percentage of the variability in systolic blood pressure is explained by age?

\vspace{2cm}

\textbf{Part F:} Use the regression equation to predict systolic blood pressure for a 50-year-old.

\vspace{3cm}

\textbf{Part G:} The residual standard error is 6.3 mmHg. What does this tell us about the model's predictions?

\vspace{3cm}

\newpage

\section{Problem 5: Comparing Methods}

A pharmaceutical company studies the relationship between drug dosage (mg) and reduction in pain score (0-10 scale) for 15 patients. Summary statistics:

\begin{itemize}
\item $\bar{X} = 50$ mg, $s_X = 15$ mg
\item $\bar{Y} = 4.2$ points, $s_Y = 1.8$ points
\item Pearson correlation: $r = 0.72$
\end{itemize}

\vspace{0.5cm}

\textbf{Part A:} Calculate the regression slope and intercept.

\vspace{0.3cm}

$b_1 = $ \rule{6cm}{0.4pt}

\vspace{2cm}

$b_0 = $ \rule{6cm}{0.4pt}

\vspace{2cm}

Regression equation: \rule{8cm}{0.4pt}

\vspace{1cm}

\textbf{Part B:} Calculate $R^2$. How much of the variability in pain reduction is explained by dosage?

\vspace{3cm}

\textbf{Part C:} Predict the pain reduction for a dosage of 60 mg.

\vspace{3cm}

\textbf{Part D:} If you wanted to achieve an average pain reduction of 5 points, what dosage would you recommend? (Hint: Solve the regression equation for X when $\hat{Y} = 5$)

\vspace{4cm}

\newpage

\section*{Reflection Questions}

Answer these questions to deepen your understanding:

\vspace{0.5cm}

\textbf{1.} What is the difference between correlation and regression? When would you use each?

\vspace{3cm}

\textbf{2.} Explain the difference between $r$ and $r^2$. Which is more directly interpretable and why?

\vspace{3cm}

\textbf{3.} Why is it important to examine a scatter plot before calculating correlation or fitting regression?

\vspace{3cm}

\textbf{4.} A study finds a strong positive correlation ($r = 0.85$) between ice cream sales and drowning deaths. Does this mean ice cream causes drowning? Explain.

\vspace{3cm}

\textbf{5.} What is extrapolation, and why is it risky in regression analysis?

\vspace{3cm}

\newpage

\section*{R Code Reference}

Here are the R commands you can use to check your work:

\vspace{0.5cm}

\textbf{For Problem 1:}
\begin{verbatim}
# Enter data
calcium <- c(800, 1000, 1200, 900, 1100, 1300, 950, 1050)
bone_density <- c(0.85, 0.90, 0.95, 0.88, 0.92, 0.98, 0.89, 0.91)

# Summary statistics
mean(calcium)
sd(calcium)
mean(bone_density)
sd(bone_density)

# Correlation test
cor.test(calcium, bone_density)

# Scatter plot
plot(calcium, bone_density,
     main = "Bone Density vs. Calcium Intake",
     xlab = "Calcium (mg)",
     ylab = "Bone Density (g/cm²)",
     pch = 19)
\end{verbatim}

\vspace{0.5cm}

\textbf{For Problems 2-3:}
\begin{verbatim}
# Fit regression model
model <- lm(bone_density ~ calcium)

# View results
summary(model)

# Extract coefficients
coef(model)

# Add regression line to plot
abline(model, col = "red", lwd = 2)

# Make predictions
predict(model, newdata = data.frame(calcium = 1100))

# Calculate residuals
residuals(model)
\end{verbatim}

\vspace{0.5cm}

\textbf{Checking assumptions:}
\begin{verbatim}
# Diagnostic plots
par(mfrow = c(2, 2))
plot(model)
par(mfrow = c(1, 1))
\end{verbatim}

\end{document}
